% ─────────────────────────────────────────────────────────────────────────────
% REVISED ABSTRACT — reviewer-proofed, calibration-honest, hedged
% ─────────────────────────────────────────────────────────────────────────────

\begin{abstract}

We describe \textsc{gcr-dosimetry-pipeline}, an open-source Python package for
estimating galactic cosmic ray (GCR) radiation exposure and cancer risk on
crewed deep-space trajectories.
The pipeline integrates Gleeson--Axford force-field solar modulation
\citep{gleeson1968} with per-species local interstellar spectra
\citep{vos2015, boschini2020},
continuous-slowing-down approximation (CSDA) slab transport with
effective-charge stopping power scaling, secondary neutron dose equivalent
from a pre-computed HZETRN table \citep{slaba2014},
and organ-specific radiation exposure in incidence of death (REID)
\citep{cucinotta2013} computed for eleven organs at ICRP-89 tissue depths.

The pipeline's per-species normalization factors are calibrated to the
MSL/RAD cruise-phase measurement of $1.84 \pm 0.33$~mGy/day
\citep{zeitlin2013} at $16\,\mathrm{g\,cm^{-2}}$ aluminum, and the
shielding-depth dependence and material ratio are independently compared
against HZETRN benchmark data \citep{slaba2014, mrigakshi2013},
showing agreement within the stated benchmark uncertainties
($\lesssim 20\%$ on absolute dose, factor of $\sim\!2$ on Al/PE ratio).
A Solar Energetic Particle module estimates BFO absorbed dose from three
canonical events using Band-function spectra \citep{tylka2006};
the August~1972 event requires $>30\,\mathrm{g\,cm^{-2}}$ aluminum
to bring BFO dose within the NASA 30-day limit of $250$~mGy.

Global physics and radiobiological uncertainties are propagated via Latin
Hypercube Sampling over eight parameters, yielding a median REID of
$\sim\!4$--$5\%$ with p5--p95 bounds spanning $\sim\!1$--$15\%$ for a
35-year-old crew member on a 259-day Earth--Mars transit at
$16\,\mathrm{g\,cm^{-2}}$ aluminum.
Variance decomposition shows that dose-and-dose-rate effectiveness factor
(DDREF) and excess risk scaling dominate the output uncertainty,
confirming that further reduction of GCR physics uncertainties alone
will have limited impact on REID bounds without commensurate advances
in high-LET human radiobiology.

Per-species dose-averaged LET decomposition reveals that carbon ions
contribute approximately $59\%$ of the unshielded GCR absorbed dose
at a field $\overline{L}_D \approx 18.7$~keV/$\mu$m, accounting for
the effective quality factor $Q_\mathrm{eff} = 2.62$.
Sixteen g/cm$^2$ of aluminum reduces the carbon fraction to $31\%$
and the field $\overline{L}_D$ to $4.1$~keV/$\mu$m,
bridging toward the LET range of clinical proton therapy beams
($0.5$--$5$~keV/$\mu$m).
The pipeline includes implementations of the Wedenberg \citep{wedenberg2013}
and McNamara \citep{mcnamara2015} proton RBE models that agree with
OpenTOPAS-RBE scorer outputs to within numerical precision for three
cell lines (\emph{in vitro} V79, prostate, head \& neck),
enabling consistent comparison of space radiation quality factors with
clinical RBE quantities from a single codebase.
All source code, 82 automated tests, validation data, and TOPAS geometry
files are freely available at \url{https://github.com/aryanhshah8/gcr-dosimetry-pipeline}.

\end{abstract}

\textbf{Keywords:}
galactic cosmic rays; Mars mission dosimetry; radiation risk; REID;
solar energetic particles; proton therapy RBE; uncertainty quantification;
open-source
